\documentclass{tutodoc}

\usepackage{../preamble.cfg}

\usepackage{../main/core-ctxt-dsl}
\usepackage{../main/main}


% TESTING LOCAL IMPLEMENTATION %

\usepackage{data-1D}


\begin{document}

\subsection{Des exemples d'utilisation directe}

\begin{tdocexa}[Une seule fonction avec les noms par défaut]
    Noter dans l'exemple suivant combien le code est similaire au rendu, ceci permettant une saisie intuitive des données.

    \tdoclatexinput{examples/data-1D/one-func-no-label.tex}
\end{tdocexa}


\begin{tdocimp}
	Retenir que tout se saisit en mode mathématique.
\end{tdocimp}


\begin{tdocwarn}
	\begin{enumerate}[wide]
		\item L'utilisation de \tdoclatexin{xvals} doit être faite une fois, et une seule, au tout début du contenu.

		\item L'emploi du nom par défaut pour une fonction n'est possible que pour une seule série d'images.
	\end{enumerate}
\end{tdocwarn}


% --------------- %


\begin{tdocexa}[Saisie rapide]
    \leavevmode

    \tdoclatexinput{examples/data-1D/shortcuts.tex}
\end{tdocexa}


% --------------- %


\begin{tdocexa}[Deux fonctions et une variable, le tout en mode \tdocquote{maison}]
    \leavevmode

    \tdoclatexinput{examples/data-1D/two-func-label.tex}
\end{tdocexa}


% --------------- %


\begin{tdocexa}[Commentaires à la sauce \LaTeX, et cellules vides]
    \leavevmode

    \tdoclatexinput{examples/data-1D/comments-empty-cells.tex}
\end{tdocexa}


% --------------- %

\begin{tdoctip}[Nombres décimaux en version \tdocquote{locale} et \tdocquote{grandes} fractions]
%    \leavevmode
%
    Via les macros \tdocmacro{dfrac} et \tdocmacro{num}%
    \footnote{
    	Cette macro ajoute de fins espaces mettant en valeur les groupes de chiffres, tout en gérant l'absence d'espaces autour de la virgule en tant que séparateur décimal.
    }
    venant des excellents packages \tdocpack{amsmath} et \tdocpack{siuntix} respectivement,
    nous pouvons rédiger des nombres décimaux, et obtenir de \tdocquote{grandes} fractions,%
    \footnote{
    	Bien que typographiquement, il ne soit pas toujours opportun d'employer de \tdocquote{grandes} fractions, ces dernières sont très utiles dans un contexte pédagogique de niveau pré-universitaire.
    }
    comme ci-dessous.

    \tdoclatexinput{examples/data-1D/num-n-dfrac.tex}
\end{tdoctip}

\end{document}
